\section{结论}
\label{sec:conclusion-zh}

本文在一个乐谱层面作曲辅助流程中整合了合法程序化语料、保留条件的 Transformer 训练、音级集合与序列分析、语法约束候选生成、MusicXML 导出和逐片段报告。在固定合成测试集上，把候选预算从 1 增至 4 并实施符号评分后，各目标自动指标相对于对齐的单候选均向优选方向变化。另一方面，从单候选模型中移除音级集合、序列、节奏或姿态字段，会使对应控制指标退化。

同一组结果也界定了方法的当前边界。神经生成没有精确实现任何完整的请求序列形式；尽管 MusicXML 结构正确，生成材料平均只覆盖约一半请求时间跨度。因此，该系统更适合被理解为可检查的草图生成器，而不是形式规则构造或作曲判断的替代品。在讨论艺术用途或程序化测试床之外的迁移之前，仍需完成盲法专家评估、独立合法乐谱验证和重复训练。
